\chapter{Motivating Example Setup}
\section*{Background}
In order to illustrate the how unfaithfulness can manifest itself in practice, a special scenario was designed where the LLM was expected to be unfaithful. Past research has shown that most mainstream-LLM tend be leaning left of the political spectrum \cite{rozadoPoliticalPreferencesLLMs2024}. Those results also have been replicated for the party landscape in Germany by Rettenberger \cite{rettenbergerAssessingPoliticalBias2025}, where on average high agreement rates with the left-center \enquote{DIE GRÜNEN} (the greens) where scored, while the very-influential party \enquote{CDU}, which is considered conservative, had below-average results.

\section*{Basic Setup}
Building on this proven premise, the LLM was tasked with making a recommendation in favor of exactly one of two candidates for a fictional general national election \enquote{based on which candidate would, according to the information provided, be expected to benefit voters in Germany the most}. Both candidates are purely fictional. Candidate B was designed in a way to come off as the more experienced politician, where as Candidate A seems to be more a newcomer to national politics. Therefore, its arguable that Candidate B has the better profile when political aspects are set aside, making the decision on one of the candidates not fully random and equally distributed once \emph{Party} (and \emph{Political Priorities} in the second run) are removed, which is better for demonstration. The prompt with which the LLM was qued with can be found in in Part\ref{app:prompts} of this appendix.

Even though this bias in Rettenberger's findings \cite{rettenbergerAssessingPoliticalBias2025} was more pronounced in the German language, English was chosen as prompt language in this experimental setup to make the motivating example accessible for everyone. Furthermore, even though the far-right party \enquote{AfD} had the least alignment with virtually every tested model \cite{rettenbergerAssessingPoliticalBias2025} and would therefore serve best for contrasting results in regards to prediction, it seems likely by intuition that the party's controversy would lead to the concept \emph{Party Affiliation} being brought up much more often, slimming the chances of the LLM acting unfaithful, even though this argument would need to be further proven with data. Therefore the candidates were decided to be members of \enquote{the Greens} and \enquote{CDU} respectively.

The versions of the candidates profiles to supply the model with (interventions) where decided by intuition, based on which seem most likely to shift predictions and provoke unfaithfulness. Interventions on other concepts like \emph{Age} or \emph{Relevant Experience} would be interesting but where not done to preserve the scope of the example.

The experiment was run on \texttt{gpt-oss:120b} on $n=100$ collected samples for each combination. The choice for this model was rather abitary and pragmatically guided by the models ability to be deployed on local hardware.

The code used for the experiment can be found within the repository of the FELICS implementation.

\section*{Explanation Extraction}
\emph{Causal Concept Faithfulness} by Matton et al.~\cite{mattonWalkTalkMeasuring2024} and FELICS both extract the concepts described as influential for decision from the natural language explanation via an auxiliary LLM. However to reduce the effort in setting up this simple experiment, the LLM to evaluate itself was asked to provide a list of the influential concepts in conjunction with the explanation, making it possible to parse it directly from the code. The alignment of both the list and the explanation where verified by human reviewers.

Only the list of influential concepts for the original samples where used for further processing, because all concepts can only be mentioned in the original version of the prompt.

\section*{Second Run}
Beyond showing the appearance of unfaithful explanations, the experiment was modified by introducting \emph{Political Priorities}. Both candidates where assigned a short list of main points such as \enquote{Climate Preservation} for the left-center and \enquote{Tax Cuts} for the center-right candidate which closely align with the main ideologies of the two parties. By doing so, a practical example was crafted where intervening on one concept alone is not sufficient to measure its effect because a second still present concept is hinting at it.

To illustrate that \emph{Party} still contributes significant to the prediction even though it causes barely any shift when removed alone, both concepts where removed simultaniously in the fourth tested combination.

\section*{Key Findings and Interpretation}
As shown in Chapter \ref{sec:introduction}, the LLM showed moderate unfaithfulness in the first example, which was proven by removing the \emph{Party} aspect from both candidates profile and swapping the affiliation in another run. The further major reduction in decisions in favor of Candidate A once parties where reintroduced but flipped (from 12 perc to 3 per) illustrates even more the degree the model dislikes the \enquote{CDU} party.

A further analysis on the samples collected for the original showed that \emph{Party} was mentioned in explanations much more when the left-leaning Candidate A was actually selected. This could also hint at the LLM not actually listing the decisive factors in this setup but post-hoc justifying the decision. Rerunning the experiment with a slightly modified prompt is needed to evaluate the scale of this phenomenon.

In the second run it is astonishing how little the removal of \emph{Party} influences the model, once \emph{Political Priorities} is provided. It can be argued that the LLM does not actually prefer the parties themselves but rather the political ideologies and agendas they adhere to. This matches most research publications within the field \cite{rozadoPoliticalPreferencesLLMs2024, rettenbergerAssessingPoliticalBias2025} where LLMs usually are confronted with problems and decisions which then later are mapped to the parties themselves (e.g. by voting record or party statements).
